% Problem record op_f408a3c300b9b214. This TeX file is derived from
% database/problems_json/op_f408a3c300b9b214.json by scripts/sync-tex.mjs; edit the JSON record
% and rerun the script rather than editing this file.
\section{Exact covariance criterion for bipartite Gaussian separability}
\label{sec:op_f408a3c300b9b214}

\paragraph{Problem.}
What necessary and sufficient condition on a covariance matrix characterizes bipartite Gaussian separability?

Let $m,n\geq1$ be integers. Let $\rho_V$ be a zero-mean Gaussian state with $m$ modes held by Alice and $n$ modes held by Bob. Its finite real covariance matrix and canonical commutators are

\begin{equation}
V_{jk}:=\operatorname{Tr}\rho_V\{R_j,R_k\},\qquad
[R_j,R_k]=i(\Omega_{m+n})_{jk},\qquad
\Omega_k:=\bigoplus_{j=1}^{k}\begin{pmatrix}0&1\\-1&0\end{pmatrix}.
\label{eq:ser16-statement-1}
\end{equation}

Separability means that $\rho_V$ belongs to the trace-norm closed convex hull of product density operators. Determine separability from the covariance in Eq.~\eqref{eq:ser16-statement-1}.

\subsection*{Status}

Solved

\subsection*{Source}

Werner and Wolf state and prove the complete covariance criterion in Proposition 1 \sourcecite{ref:ser16-1}{WW01}.

\subsection*{Progress}

\begin{itemize}
  \item Werner--Wolf, Proposition 1, proves both directions of the criterion below. The exact criterion is the following finite semidefinite feasibility problem over real symmetric matrices $V_A$ of size $2m$ and $V_B$ of size $2n$:

\begin{equation}
\rho_V\text{ is separable}
\quad\Longleftrightarrow\quad
\exists V_A,V_B:\quad V_A+i\Omega_m\geq0,\quad
V_B+i\Omega_n\geq0,\quad V\geq V_A\oplus V_B.
\label{eq:ser16-progress-1-1}
\end{equation}

Equation~\eqref{eq:ser16-progress-1-1} solves the unrestricted Gaussian separability criterion. \sourcecite{ref:ser16-1}{WW01}

  \item With $T_B:=I_{2m}\oplus\bigoplus_{j=1}^{n}\operatorname{diag}(1,-1)$, partial-transpose positivity is equivalent to

\begin{equation}
T_BVT_B+i\Omega_{m+n}\geq0.
\label{eq:ser16-progress-2-1}
\end{equation}

Equation~\eqref{eq:ser16-progress-2-1} is also sufficient for separability if $m=1$ or $n=1$, or if the Gaussian state is invariant under all permutations of the modes on one party; it fails to be sufficient for general $m=n=2$. \sourcecite{ref:ser16-1}{WW01}, \sourcecite{ref:ser16-2}{LSA18} The permutation-symmetric extension is Theorem 9 of Lami et al. \sourcecite{ref:ser16-2}{LSA18}
\end{itemize}

\subsection*{References}

\begin{enumerate}[label={},leftmargin=4.8em,labelsep=0.5em]
  \footnotesize
  \item[\textup{[WW01]}]\label{ref:ser16-1}
  R. F. Werner and M. M. Wolf, "Bound Entangled Gaussian States," \emph{Physical Review Letters} \textbf{86}, 3658--3661 (2001). \href{https://doi.org/10.1103/PhysRevLett.86.3658}{doi:10.1103/PhysRevLett.86.3658}; \href{https://arxiv.org/abs/quant-ph/0009118}{arXiv:quant-ph/0009118}.

  \item[\textup{[LSA18]}]\label{ref:ser16-2}
  L. Lami, A. Serafini, and G. Adesso, "Gaussian Entanglement Revisited," \emph{New Journal of Physics} \textbf{20}, 023030 (2018). \href{https://doi.org/10.1088/1367-2630/aaa654}{doi:10.1088/1367-2630/aaa654}; \href{https://arxiv.org/abs/1612.05215}{arXiv:1612.05215}.
\end{enumerate}

\subsection*{Comment}

The full criterion is solved by peer-reviewed results. Semidefinite feasibility is an exact mathematical characterization. A demand for a particular elementary expression would require a separately specified expression class.

\subsection*{Field}

Quantum Resource Theory

\subsection*{Topic}

Gaussian quantum information; Quantum separability

\subsection*{ID}

\texttt{op\_f408a3c300b9b214}
